\documentclass[11pt]{amsart}
\usepackage[T1]{fontenc}
\usepackage{lmodern}
\usepackage{amsmath,amssymb,amsthm,mathtools}
\usepackage{xcolor}
\usepackage[colorlinks=true,linkcolor=blue!55!black,citecolor=blue!55!black,urlcolor=blue!55!black]{hyperref}

\newtheorem{theorem}{Theorem}[section]
\newtheorem{lemma}[theorem]{Lemma}
\newtheorem{proposition}[theorem]{Proposition}
\newtheorem{corollary}[theorem]{Corollary}
\theoremstyle{definition}
\newtheorem{definition}[theorem]{Definition}
\newtheorem{remark}[theorem]{Remark}

\newcommand{\RV}{\mathrm{RV}}
\newcommand{\TC}{\mathrm{TC}}
\newcommand{\D}{\mathcal D}

\title[The $k=4$ phylogenetic-network conjecture is false]
{A square-root obstruction to the $k=4$ phylogenetic-network enumeration conjecture}
\author{Zijian Zeng}
\address{Institute of Computer Science and Digital Innovation, UCSI University,
Kuala Lumpur 56000, Malaysia}
\email{zijianzeng@foxmail.com}
\email{1002266693@ucsiuniversity.edu.my}
\date{August 29, 2026}
\subjclass[2020]{05A15, 05C30, 92B10}
\keywords{phylogenetic network, reticulation-visible network, tree-child network,
component graph, exponential generating function}
\hypersetup{
  pdftitle={A square-root obstruction to the k=4 phylogenetic-network enumeration conjecture},
  pdfauthor={Zijian Zeng}
}

\begin{document}

\begin{abstract}
Batle conjectured that, for fixed reticulation number $k$, the difference
between the numbers of reticulation-visible and tree-child networks is a sum
of two prescribed factorial bases with polynomial coefficients.  We disprove
the conjecture at its first open case, $k=4$.  If
$X=(1-2z)^{1/2}$ and $R_4(z)$ is the exponential generating function for
reticulation-visible networks, the Chang--Fuchs component-graph formula gives
$[X^3]R_4(z)=5/16$.  A highest-exponent argument isolates this coefficient in
the unique four-sink component graph, so the disproof does not depend on a
large numerical enumeration.  Neither proposed factorial basis can produce
a positive odd power $X^j$ with $j>1$.  We also compute the corrected exact
formula: it contains the additional term $15(2\ell-5)!!/16$.  Two independent
exact implementations certify the full Laurent generating function and an
integer residual against the conjectured degree bounds.
\end{abstract}

\maketitle

\section{Introduction}

A binary phylogenetic network is a rooted directed acyclic graph whose leaves
are labelled and whose internal vertices are tree vertices or reticulation
vertices.  Two important subclasses are the tree-child networks and the
reticulation-visible networks.  Chang and Fuchs developed a component-graph
method for exact enumeration with a fixed small number of reticulations
\cite{ChangFuchs2024}.  The tree-child counts are connected to the
Pons--Batle word numbers \cite{PonsBatle2021}; the required identity has since
been proved in full generality \cite{LinEtAl2026}, with a further combinatorial
framework in \cite{LiuWallnerYu2026}.

For the sets $\RV_{\ell,k}$ and $\TC_{\ell,k}$ of networks with $\ell$ leaves
and $k$ reticulations, put
\[
 \Delta_k(\ell)=|\RV_{\ell,k}|-|\TC_{\ell,k}|.
\]
Batle's Conjecture~5.1 \cite{Batle2026} asserts, for $\ell\ge k+1$, that
\begin{equation}\label{eq:conjecture}
 \Delta_k(\ell)=A_k(\ell)(2\ell-3)!!
 -B_k(\ell)2^{\ell-k+1}\ell!,
\end{equation}
where $A_k,B_k$ are polynomials of degrees $2k-1$ and $2k-3$, respectively,
and the leading coefficient of $A_k$ is $2^k$.  The cases $k=2,3$ motivated
the conjecture; the case $k=4$ was explicitly left open.

Our main result is stronger than the negation of the degree conditions.

\begin{theorem}\label{thm:main}
There do not exist polynomials $A,B\in\mathbb Q[x]$ of any finite degrees such
that
\[
 \Delta_4(\ell)=A(\ell)(2\ell-3)!!-B(\ell)2^{\ell-3}\ell!
\]
for every integer $\ell\ge5$.  In particular, the $k=4$ specialization of
\eqref{eq:conjecture} is false.
\end{theorem}

In fact the conjecture misses one lower double-factorial term.  Define
\begin{align*}
A_0(x)={}&8x^7+\frac{110}{3}x^6-\frac{1027}{105}x^5
-\frac{34739}{210}x^4-\frac{1088}{5}x^3\\
&-\frac{7831}{140}x^2+\frac{48467}{840}x+\frac{309}{16},\\
B_0(x)={}&64x^5+82x^4-\frac{985}{6}x^3
-\frac{1091}{2}x^2-\frac{1633}{3}x-194.
\end{align*}

\begin{theorem}\label{thm:corrected}
For every integer $\ell\ge5$,
\begin{equation}\label{eq:corrected}
 \Delta_4(\ell)=A_0(\ell)(2\ell-3)!!
 +\frac{15}{16}(2\ell-5)!!
 -B_0(\ell)2^{\ell-3}\ell!.
\end{equation}
\end{theorem}

Thus the coefficient of $(2\ell-3)!!$ in the corrected formula is
$A_0(\ell)+15/(16(2\ell-3))$, not a polynomial.  The proof of
Theorem~\ref{thm:main} is human-checkable and uses only one component graph.
The full exact sum is needed for Theorem~\ref{thm:corrected}; it is accompanied
by two independent rational-arithmetic implementations.

The literature search for this note used exact-title, conjecture, formula,
citation, and correction queries through August 29, 2026.  No other solution
of the $k=4$ case was found in those searches.  This is a bounded search
statement, not a claim about inaccessible or unindexed work.

\section{Definitions and the component-graph formula}

A binary phylogenetic network on $[\ell]=\{1,\dots,\ell\}$ is a rooted simple
directed acyclic graph.  Its root has indegree zero and outdegree one; its
leaves have indegree one and outdegree zero and are bijectively labelled by
$[\ell]$; every other vertex has degree pair $(1,2)$ or $(2,1)$.  The latter
vertices are reticulations.  Isomorphisms fix all leaf labels.

A network is tree-child if every non-leaf vertex has a tree vertex or a leaf
as a child.  A reticulation is visible if there is a leaf such that every path
from the root to that leaf contains the reticulation.  A network is
reticulation-visible if all its reticulations are visible.  These are the
conventions in \cite{Batle2026,ChangFuchs2024}.

Let $M_{\ell,d}$ count one-component networks in the normalization of
Chang--Fuchs, and define
\[
 F_d(z)=\sum_{r\ge0}M_{r+d,d}\frac{z^r}{r!}.
\]
Put $X=(1-2z)^{1/2}$.  In particular, $F_0(z)=1-X$, with $M_{0,0}=0$.

Let $\D_m$ be the set of unlabelled rooted directed acyclic multigraphs with
$m$ vertices, every non-root vertex having indegree two, and with double
edges allowed.  For $G\in\D_m$ and $v\in V(G)$, let $c(v)$ be the number of
distinct children and $d(v)$ the number of children joined to $v$ by a double
edge.  Let $m(G)$ be the order of the root-preserving automorphism group.
Chang--Fuchs Proposition~12 gives
\begin{equation}\label{eq:CF}
 R_k(z):=\sum_{\ell\ge1}|\RV_{\ell,k}|\frac{z^\ell}{\ell!}
 =\sum_{G\in\D_{k+1}}\frac1{m(G)}
   \prod_{v\in V(G)}W_{c(v),d(v)}(z).
\end{equation}
Directly shifting the index in their vertex series gives
\begin{equation}\label{eq:weights}
 W_{c,d}=F_d^{(c-d)}\quad(d>0),\qquad
 W_{c,0}=F_0^{(c)}-M_{c,0}.
\end{equation}
Here differentiation in $z$ acts on Laurent polynomials in $X$ as
$-X^{-1}d/dX$.

\section{The isolated square-root obstruction}

The one-component recurrence of Gunawan, Rathin, and Zhang
\cite{GunawanRathinZhang2020}, as recorded in \cite{ChangFuchs2024}, is
\begin{align}\label{eq:Mrec}
M_{\ell,d}={}&(\ell+d-2)M_{\ell,d-1}+(d-1)M_{\ell,d-2}\notag\\
&+\frac12\sum_{j=1}^{d-1}\binom{d-1}{j}(2j-1)!!
 \bigl(M_{\ell-j,d-1-j}-M_{\ell+1-j,d-1-j}\bigr),
\end{align}
with $M_{\ell,0}=(2\ell-3)!!$, $M_{\ell,1}=(\ell-1)(2\ell-3)!!$, and
$M_{\ell,d}=0$ for $\ell<d$.  Chang--Fuchs Lemma~7 implies
\begin{equation}\label{eq:Fdshape}
 F_d(z)=\frac{N_d(z)}{X^{4d-1}},\qquad \deg N_d\le2d-1.
\end{equation}

For $d=4$, the eight values from \eqref{eq:Mrec} are
\begin{align*}
&6249,\ 97182,\ 1701855,\ 33121890,\\
&709428825,\ 16587636030,\ 420498508815,\ 11489311483650.
\end{align*}
Multiplying $F_4$ by $(1-2z)^{15/2}$, or equivalently taking the finite
Cauchy product, uses
\[
 [z^r]N_4(z)=\sum_{m=0}^r\frac{M_{4+r-m,4}}{(r-m)!}
 \binom{15/2}{m}(-2)^m\qquad(0\le r\le7)
\]
and gives
\begin{align}\label{eq:N4}
N_4(z)={}&6249+3447z+2475z^2-2370z^3+4215z^4\notag\\
&-4968z^5+3216z^6-960z^7.
\end{align}
All quantities in this section are integers or rationals; no numerical
approximation is involved.

For a nonzero Laurent polynomial $H$, write $h(H)$ for the largest exponent
of $X$ occurring in $H$.

\begin{lemma}\label{lem:highest}
Among the summands in \eqref{eq:CF} for $k=4$, only the graph consisting of
the root joined by a double edge to each of four sinks can contain $X^3$.
Its $X^3$ coefficient is $5/16$.
\end{lemma}

\begin{proof}
By \eqref{eq:Fdshape}, $h(F_d)\le-1$ for $d>0$.  Each differentiation in $z$
lowers every surviving exponent by two, so \eqref{eq:weights} gives
\[
 h(W_{c,d})\le-1-2(c-d)\quad(d>0).
\]
For $d=0$, $h(W_{0,0})=1$, whereas $h(W_{c,0})\le0$ for $c>0$.

Suppose $G\in\D_5$ has $s$ sinks and $q$ vertices with at least one outgoing
double edge.  It follows that
\begin{equation}\label{eq:heightbound}
 h\left(\prod_vW_{c(v),d(v)}\right)\le s-q.
\end{equation}
If $s\le2$, the right side is at most two.  If $s=3$, the two nonsinks are
the root and one non-root vertex.  The latter has indegree two; if there were
no double edge, it would need two distinct nonsink ancestors although only the
root is available.  Hence $q\ge1$, and \eqref{eq:heightbound} is again at most
two.

If $s=4$, the only nonsink is the root, so every sink is joined to it by a
double edge.  This graph is unique and has automorphism order $4!=24$.  Its
contribution is $F_4(z)(1-X)^4/24$.  To obtain $X^3$, the $X^{14}$ term of
$N_4((1-X^2)/2)$ must be multiplied by the $X^4$ term of $(1-X)^4$.  By
\eqref{eq:N4}, the former is
\[
 -960\left(-\frac{X^2}{2}\right)^7=\frac{15}{2}X^{14}.
\]
Division by $24$ gives $5/16$.
\end{proof}

\begin{corollary}\label{cor:x3}
The exact generating function satisfies $[X^3]R_4(z)=5/16$.
\end{corollary}

\section{The proposed bases cannot contain the obstruction}

Let $\theta=z\,d/dz$.  Since
\[
 \sum_{\ell\ge1}(2\ell-3)!!\frac{z^\ell}{\ell!}=1-X,
\]
the exponential generating function of
$P(\ell)(2\ell-3)!!$ is $P(\theta)(1-X)$.  On a Laurent monomial,
\begin{equation}\label{eq:theta}
 \theta X^e=\frac e2\bigl(X^e-X^{e-2}\bigr).
\end{equation}
Starting from $X$, repeated use of \eqref{eq:theta} produces only odd powers
whose exponent is at most one.  Thus this basis cannot produce $X^3$.

For the other basis,
\[
 \sum_{\ell\ge0}Q(\ell)2^{\ell-3}\ell!\frac{z^\ell}{\ell!}
 =\frac18\sum_{\ell\ge0}Q(\ell)(2z)^\ell.
\]
This is rational in $2z=1-X^2$, with denominator a power of $X^2$, and hence
has only even powers of $X$.  Changing finitely many initial coefficients adds
a polynomial in $z=(1-X^2)/2$, which again has only even powers.  We have
proved the following span obstruction.

\begin{lemma}\label{lem:span}
For arbitrary finite-degree polynomials $P,Q$, the eventual exponential
generating function of
\[
 P(\ell)(2\ell-3)!!+Q(\ell)2^{\ell-3}\ell!
\]
has no $X^j$ term for any positive odd $j>1$.
\end{lemma}

The proved Pons--Batle identity and their equation~(19d) give
\begin{align}\label{eq:TC4}
|\TC_{\ell,4}|=\binom\ell4\bigl[&(2\ell+5)!!-4(2\ell+4)!!
+2(2\ell+3)!!\notag\\
&+\tfrac{17}{2}(2\ell+2)!!-\tfrac{283}{63}(2\ell+1)!!\bigr]
\end{align}
for $\ell\ge5$ \cite{PonsBatle2021,LinEtAl2026}.  Expressing the shifted odd
double factorials relative to $(2\ell-3)!!$ and the even double factorials
relative to $2^{\ell-3}\ell!$ shows that \eqref{eq:TC4} lies in the span of
Lemma~\ref{lem:span}.

\begin{proof}[Proof of Theorem~\ref{thm:main}]
If the asserted representation of $\Delta_4$ existed, adding
\eqref{eq:TC4} would put $|\RV_{\ell,4}|$ in the same eventual two-basis span.
Lemma~\ref{lem:span} would force $[X^3]R_4=0$, contradicting
Corollary~\ref{cor:x3}.
\end{proof}

\section{The corrected closed form}

For completeness, exact evaluation of all component graphs in $\D_5$ gives
\begin{align}\label{eq:R4full}
R_4(z)={}&\frac5{16}X^3-\frac54X^2+\frac{97}{32}X-4
+\frac{45}{32}X^{-1}+\frac{13}{8}X^{-2}-\frac{331}{128}X^{-3}\notag\\
&-\frac{17}{8}X^{-4}+\frac{1111}{64}X^{-5}-\frac{19}{8}X^{-6}
+\frac{4487}{128}X^{-7}+\frac{385}{8}X^{-8}\notag\\
&-\frac{319}{4}X^{-9}-586X^{-10}+\frac{45835}{128}X^{-11}
+1045X^{-12}-\frac{45045}{64}X^{-13}\notag\\
&-480X^{-14}+\frac{45045}{128}X^{-15}.
\end{align}
This identity was computed twice by the independent procedures described in
Section~\ref{sec:certificate}.

For $j\ge1$ and sufficiently large $\ell$, the binomial theorem gives
\begin{align*}
\ell![z^\ell]X^{-(2j-1)}
 &=\frac{(2\ell+2j-3)!!}{(2j-3)!!},\\
\ell![z^\ell]X^{-2j}
 &=2^\ell\ell!\binom{\ell+j-1}{j-1}.
\end{align*}
Moreover,
\[
\ell![z^\ell]X=-(2\ell-3)!!,
\qquad
\ell![z^\ell]X^3=3(2\ell-5)!!.
\]
The nonnegative even powers in \eqref{eq:R4full} have zero coefficient for
$\ell\ge5$.  Substituting these identities into \eqref{eq:R4full}, subtracting
\eqref{eq:TC4}, and collecting terms gives exactly \eqref{eq:corrected}.
This proves Theorem~\ref{thm:corrected}.

\section{Exact certificate and independent checks}\label{sec:certificate}

The discovery implementation uses SymPy rational arithmetic.  With the root
fixed, it chooses the unordered pair of parents of each of the four non-root
vertices, rejects precisely the cyclic choices, evaluates the vertex weights,
and divides the resulting labelled sum by $4!$.

The independent verifier uses only the Python standard library.  It generates
the graphs through all topological orders, deduplicates them, canonically
quotients by root-preserving permutations, and weights every class by the
reciprocal of its independently computed automorphism order.  It implements
Laurent arithmetic with \texttt{fractions.Fraction}.  Both routes give 1587
fixed-root labelled graphs and 79 unlabelled classes, with automorphism-order
histogram
\[
56[1]+19[2]+[4]+2[6]+[24].
\]
A third check uses the sink-row inclusion--exclusion recurrence and gives 7935
fully vertex-labelled graphs, equal to $5\cdot1587$.

Cardona and Zhang's Figure~9 is captioned as a list of 82 five-vertex classes
\cite{CardonaZhang2020}.  The 82 orbit sizes printed in that figure sum to
7905, not 7935.  Thus that figure is not used as input here; exact generation
from the definition and the independent recurrence are used instead.  This
discrepancy is immaterial to Theorem~\ref{thm:main}, whose proof isolates the
unique four-sink graph without enumerating the other classes.

As a further check, the 14 coefficients allowed by
$\deg A\le7$, $\deg B\le5$ are uniquely fitted from exact values at
$\ell=5,\dots,18$.  At $\ell=19$ the actual value is
\[
618911914073871797682864561000,
\]
whereas the fit predicts
\[
618911914073871699607286961000.
\]
The exact residual is $98075577600000>0$.  This finite computation is a
diagnostic falsification of the stated degree bounds; it is not used as proof
of the all-$\ell$ assertion.

The accompanying archive contains the stable paths
\begin{verbatim}
code/discover_delta4_sympy.py
code/verify_delta4_certificate.py
certificates/delta4_laurent.json
\end{verbatim}
and the full verification command is
\begin{verbatim}
python3 code/verify_delta4_certificate.py
\end{verbatim}
It checks the isolated-star coefficient, the complete Laurent polynomial, the
three graph counts, the corrected formula on $5\le\ell\le30$, and the
$\ell=19$ residual using exact arithmetic.  The finite-range check is redundant
evidence for Theorem~\ref{thm:corrected}, whose proof is the coefficient
extraction above.

\section{Scope and disclosure}

The result addresses only Batle's conjectured factorial-basis form and its
$k=4$ case.  It does not make a claim about the conjecture at every larger
$k$, nor does it alter the definitions or the Chang--Fuchs structural theorem.
The main proof depends on their Proposition~12 and Lemma~7, and the subtraction
step depends on the now-proved Pons--Batle identity.

Every decisive computation uses rational or integer arithmetic and has an
independent implementation.  No Lean, Coq, Isabelle, or other proof assistant
was used.  The manuscript has not undergone journal peer review.

\bibliographystyle{amsplain}
\begingroup
\sloppy
\bibliography{references}
\endgroup

\end{document}
